\section{Discussion and limitations}
\label{sec:discussion}

A shared encoder with SIGReg can support a Gaussian residual predictor on Moving-MNIST. Relative to \shdet{}, the gain combines forecast spread with a more decodable target representation. Relative to the \emavar{}, the observed mean scores are close, but five replications leave the difference uncertain. On MPI3D, none of the ten variational-model runs beats its own persistence forecast on the latent score, and the deterministic models cannot by construction.

The four-model design changes several components together. EMA versus Shared alters target gradients, regularization and BN exposure; deterministic versus Gaussian residual prediction adds capacity and posterior information during training. We cannot attribute the results to randomness, SIGReg or the removal of EMA alone. The primary metric includes a fitted linear readout, and its distortion contributes to \shdet{}'s Moving-MNIST error and to the MPI3D scores. Other readouts or prediction normalization could change the results. The two tasks also use different architectures and posterior inputs, so differences between them cannot be attributed to the data alone.

Five replications describe training randomness on fixed splits and evaluation banks, and the analysis includes no equivalence test. Both tasks are synthetic, with single-step, low-dimensional targets. Each dataset uses one optimization recipe for all four models, and we did not retune it for the Shared models. Equal update budgets need not give equal convergence: we could not inspect training curves, and on MPI3D the learning rate is still about 11\% of its peak at the final update. We did not test sensitivity to $\beta$, $\lambda$ or the SIGReg settings, or establish Monte Carlo and quadrature convergence of the scores.

On Moving-MNIST, the \emavar{} follows the published AdaSSL video recipe with documented local conventions, because the official repository, at the commit we pinned, contains no video pipeline (Appendix~\ref{app:fidelity}). On MPI3D, both EMA references are our own image-and-command adaptations, so their failure there should not be attributed to the published methods.

The MPI3D split was inspected during earlier development experiments that are not reported here, and earlier use of the MNIST test set cannot be ruled out, although the development protocol recorded no test access. Checkpoints, fitted readouts, per-query outputs and training curves were not available for this analysis, so it uses exported run-level results only and shows no forecast examples or training curves. We checked the method descriptions against the experiment code and configuration records, but could not verify byte for byte that this code is the snapshot that was executed.
